% Semantic Hierarchy
\documentclass[tikz,border=10pt]{standalone}
\usepackage[robusto]{diagrama}
\begin{document}
\begin{tikzpicture}[img, font=\sffamily]
\tikzset{
  ov/.style={line width=1.2pt},
  t/.style={font=\sffamily\px{10.5}},
  tb/.style={font=\sffamily\bfseries\px{11}},
  rama/.style={line width=2.6pt},
  lista/.style={font=\sffamily\px{8.2},anchor=north west, align=left, inner sep=0pt},
  llavep/.style={puntitos, line width=1.2pt, llave=0.3cm},
}
% bifurcaciones: « < » desde un óvalo y « > » hacia el siguiente
\newcommand\abre[1]{\draw[rama] (#1+0.25,1.21) -- (#1+0.17,1.21) -- (#1,1.38) -- (#1+0.17,1.59) -- (#1+0.25,1.59);}
\newcommand\cierra[1]{\draw[rama] (#1-0.30,1.21) -- (#1-0.20,1.21) -- (#1,1.38) -- (#1-0.20,1.59) -- (#1-0.30,1.59);}

\node[tb] at (3.55,0.59) {Semantic Hierarchy};
\draw[puntitos, line width=1.4pt, -{Straight Barb[length=2.4mm,width=3.2mm]}] (0.28,0.80) -- (7.00,0.80);

\draw[ov] (0.48,1.38) ellipse (0.40 and 0.25);
\node[tb, font=\sffamily\bfseries\px{9.5}] at (0.48,1.38) {RDRD /\\Classification};
\abre{0.88} \cierra{2.17}
\node[t] at (1.52,1.13) {Non-\\Conceptual};
\node[t] at (1.52,1.62) {Proto-\\Conceptual\\(Labelling)};

\draw[ov] (2.48,1.38) ellipse (0.31 and 0.29);
\node[tb, font=\sffamily\bfseries\px{9.5}] at (2.48,1.38) {Description\\(Material\\inference)};
\abre{2.80} \cierra{4.17}
\node[t] at (3.56,1.22) {Incompatibility};
\node[t] at (3.56,1.59) {Consequence};

\draw[ov] (4.60,1.37) ellipse (0.43 and 0.27);
\node[tb, font=\sffamily\bfseries\px{9.5}] at (4.60,1.37) {Logical\\explicitation\\(Formal\\inference)};
\abre{5.04} \cierra{6.35}
\node[t] at (5.70,1.17) {Pragmatic\\Inferences};
\node[t] at (5.70,1.60) {Semantic\\Inferences};

\draw[ov] (6.71,1.40) ellipse (0.36 and 0.26);
\node[tb] at (6.71,1.40) {Theory\\formation};
\draw[rama] (7.13,1.24) -- (7.73,0.20) -- (7.78,0.20);
\draw[rama] (7.12,1.44) -- (7.53,1.83) -- (7.75,1.83);

% teorías
\node[tb, anchor=north west, align=left] at (7.82,0.04) {Informal\\Theories};
\node[lista] at (7.88,0.33) {\textbullet\ Hermeneutics\\\textbullet\ Meta-ontologies};
\node[tb, anchor=north west] at (7.78,0.86) {Formal Theories};
\node[lista] at (7.82,1.04) {%
  \textbullet\ S1) Data Recollection as structure identification\\
  \phantom{\textbullet\ }for informal theories.\\
  \textbullet\ S2) Formalization of theory\\
  \quad\textbullet\ Coherence-Network Method\\
  \qquad\textbullet\ Non-linear systems\\
  \quad\textbullet\ Axiomatics\\
  \qquad\textbullet\ Euclidean\\
  \qquad\textbullet\ Informal Hilbertean\\
  \qquad\textbullet\ Formal Hilbert\\
  \qquad\textbullet\ Informal Set theory\\
  \qquad\textbullet\ Formal set theory\\
  \textbullet\ S3) Systematization of theories - Meta theories\\
  \quad\textbullet\ Partial (Systematization of some theories)\\
  \quad\textbullet\ Global systematics (of the universe of discourse)\\
  \textbullet\ S4) Assessment of system.};
\draw[puntitos, line width=1.3pt, {Straight Barb[length=2.4mm,width=3.2mm]}-{Straight Barb[length=2.4mm,width=3.2mm]}]
  (8.97,1.36) -- (9.83,1.36) -- (9.83,2.34) -- (9.40,2.34);

% monotonicidad
\draw[puntitos, line width=1.4pt, -{Straight Barb[length=2.4mm,width=3.2mm]}] (9.95,0.00) -- (9.95,3.02);
\node[font=\sffamily\px{13}, rotate=-90] at (10.30,1.40) {Monotonicity};

% reglas y normas
\draw[llavep] (7.05,1.75) -- (0.20,1.75);
\node[tb, font=\sffamily\bfseries\px{12}] at (3.53,2.11) {Obeying Rules};
\draw[llavep] (4.25,2.13) -- (0.22,2.13);
\node[tb, font=\sffamily\bfseries\px{12}] at (2.30,2.49) {Following Norms};
\draw[llavep] (2.22,2.62) -- (0.10,2.62);
\node[tb, font=\sffamily\bfseries\px{12}] at (1.18,2.99) {Accordance to Rules};
\end{tikzpicture}
\end{document}
